% Old stuff regarding soundness, will eventually delete.








\subsection{Soundness (NEW)}

This section states the intended soundness argument using the symbolic
semantics as the main coupling object.  A symbolic state has the form
\[
  \langle \Sigma \mid\mid e \rangle,
\]
where $\Sigma$ is a sequence of symbolic samples
\[
  u_1 \sim D_1;\ldots;u_k \sim D_k.
\]
The arguments to $D_i$ are affine expressions over earlier symbolic variables
$u_1,\ldots,u_{i-1}$.

Let $\mathsf{actual}$ and $\mathsf{expected}$ be the two interpretations of
symbolic configurations.  The actual interpretation samples the symbolic
environment and then returns the residual expression:
\[
  \mathsf{actual}(\langle \Sigma \mid\mid e \rangle)
  =
  u_1 \sim D_1;\ldots;u_k \sim D_k;\mathsf{return}\ e.
\]
The expected interpretation substitutes the expectation of each symbolic
sample, in order, into the later symbolic environment and residual expression:
\[
  \mathsf{expected}(\langle \Sigma \mid\mid e \rangle)
  =
  \mathsf{return}
  e[
    \mathbb{E}[D_1]/u_1,\ldots,
    \mathbb{E}[D_k]/u_k
  ].
\]
Both interpretations extend from symbolic states to symbolic configurations by
bind.

\begin{definition}[Type-indexed observational equivalence]
  We write $\mu_1 \approx_\tau \mu_2$ for the equivalence relation induced by
  type $\tau$.  At expectation-mode floats this means equal expectations,
  while at general-mode floats, booleans, and the other non-expectation
  observations it means equality of distributions:
  \[
    \mu_1 \approx_{\Float{\E}} \mu_2
    \quad\Longleftrightarrow\quad
    \mathbb{E}[\mu_1] = \mathbb{E}[\mu_2],
  \]
  and
  \[
    \mu_1 \approx_{\Float{\G}} \mu_2
    \quad\Longleftrightarrow\quad
    \mu_1 = \mu_2.
  \]
  and
  \[ 
    \mu_1 \approx_{\Bool} \mu_2
    \quad\Longleftrightarrow\quad
    \mu_1 = \mu_2.
  \]

  The relation is extended structurally to products, sums, lists, and
  functions.
\end{definition}

\begin{lemma}[Determinization preservation]
  If $\emptyset \vdash e : \tau$, then
  \[
    \emptyset \vdash \exptrans{e : \tau} : \tau.
  \]
\end{lemma}

\begin{lemma}[Symbolic preservation]
  If a symbolic configuration $C$ is well typed at type $\tau$, and
  \[
    C \longrightarrow_{\sym} C',
  \]
  then $C'$ is also well typed at type $\tau$.  Consequently, if
  \[
    \langle \emptyset \mid\mid e \rangle
    \longrightarrow_{\sym}^{n}
    C_n
  \]
  and $\emptyset \vdash e : \tau$, then $C_n$ is well typed at type
  $\tau$.
\end{lemma}

\begin{lemma}[Actual interpretation]
  If
  \[
    \langle \emptyset \mid\mid e \rangle
    \longrightarrow_{\sym}^{n}
    C_n,
  \]
  then
  \[
    \mathsf{actual}(C_n) = \bigsemstep{e : \tau}{n}.
  \]
\end{lemma}

\begin{lemma}[Expected interpretation]
  If $\emptyset \vdash e : \tau$ and
  \[
    \langle \emptyset \mid\mid e \rangle
    \longrightarrow_{\sym}^{n}
    C_n,
  \]
  then
  \[
    \mathsf{expected}(C_n)
    \approx_\tau
    \bigsemstep{\exptrans{e : \tau} : \tau}{n}.
  \]
\end{lemma}

\begin{lemma}[Well-typed symbolic configurations are expectation-affine]
  If $C$ is a well-typed symbolic configuration at type $\tau$, then its
  residual expression is expectation-affine in every expectation-mode symbolic
  variable.  Therefore
  \[
    \mathsf{actual}(C) \approx_\tau \mathsf{expected}(C).
  \]
\end{lemma}

\begin{theorem}[$n$-step symbolic soundness]
  If $\emptyset \vdash e : \tau$, then for every $n \in \mathbb{N}$,
  \[
    \bigsemstep{e : \tau}{n}
    \approx_\tau
    \bigsemstep{\exptrans{e : \tau} : \tau}{n}.
  \]
\end{theorem}

\begin{proof}
  Start with a closed well-typed source expression
  \[
    \emptyset \vdash e : \tau.
  \]
  By determinization preservation,
  \[
    \emptyset \vdash \exptrans{e : \tau} : \tau.
  \]
  Embed the source expression into the symbolic semantics and run it for
  $n$ steps:
  \[
    \langle \emptyset \mid\mid e \rangle
    \longrightarrow_{\sym}^{n}
    C_n.
  \]
  By symbolic preservation, $C_n$ is well typed at type $\tau$.  The actual
  interpretation theorem gives
  \[
    \mathsf{actual}(C_n) = \bigsemstep{e : \tau}{n},
  \]
  and the expected interpretation theorem gives
  \[
    \mathsf{expected}(C_n)
    \approx_\tau
    \bigsemstep{\exptrans{e : \tau} : \tau}{n}.
  \]
  Since $C_n$ is well typed, it is expectation-affine, so
  \[
    \mathsf{actual}(C_n) \approx_\tau \mathsf{expected}(C_n).
  \]
  Therefore, by transitivity of $\approx_\tau$,
  \[
    \bigsemstep{e : \tau}{n}
    \approx_\tau
    \bigsemstep{\exptrans{e : \tau} : \tau}{n}.
  \]
\end{proof}

\begin{theorem}[Soundness of determinization]
  If $\emptyset \vdash e : \tau$, then
  \[
    \bigsem{e : \tau}
    \approx_\tau
    \bigsem{\exptrans{e : \tau} : \tau}.
  \]
  In particular, if $\tau = \Float{\E}$, then
  \[
    \mathbb{E}[\bigsem{e : \Float{\E}}]
    =
    \mathbb{E}[\bigsem{\exptrans{e : \Float{\E}} : \Float{\E}}].
  \]
\end{theorem}

\begin{proof}
  The $n$-step symbolic soundness theorem gives
  \[
    \bigsemstep{e : \tau}{n}
    \approx_\tau
    \bigsemstep{\exptrans{e : \tau} : \tau}{n}
  \]
  for every $n$.  Passing to the big-step semantics gives
  \[
    \bigsem{e : \tau}
    \approx_\tau
    \bigsem{\exptrans{e : \tau} : \tau}.
  \]
  The expectation-mode float case follows by the definition of
  $\approx_{\Float{\E}}$.
\end{proof}

\subsection{Soundness}
\begin{definition}[Expected value]
  Suppose $\mu$ is a probability measure on $\mathbb{R}$. Then, the expected value of the distribution $\mu$ is:
  \[
    \mathbb{E}[\mu] = \int_{\mathbb{R}} x d\mu(x).
  \]
\end{definition}

\begin{lemma}[One-step value]\label{lem:one-step-value}
  If $e : \tau \in Val_{\tau}$, then 
  \[
    \sem{e : \tau} = \delta_{e : \tau}.
  \]
\end{lemma}

\begin{proof}
  By~\Cref{tab:small-step-semantics}.
\end{proof}

\begin{lemma}[$n$-step value]\label{lem:nstep-value}
  If $e : \tau \in Val_{\tau}$, then
  \[
    \bigsemstep{e : \tau}{n} = \delta_{e : \tau}
  \]
  for every $n \in \mathbb{N}$.
\end{lemma}

\begin{proof}
  By induction on $n$.

  \textbf{Base case.}
  Suppose $n = 0$. By definition of the $0$-step semantics,
  \[
    \bigsemstep{e : \tau}{0}
    =
    \delta_{e : \tau}.
  \]

  \textbf{Inductive step.}
  Suppose
  \[
    \bigsemstep{e : \tau}{n}
    =
    \delta_{e : \tau}.
  \]
  We must show
  \[
    \bigsemstep{e : \tau}{n+1}
    =
    \delta_{e : \tau}.
  \]
  By definition of the $(n+1)$-step semantics,
  \[
    \bigsemstep{e : \tau}{n+1}
    =
    \bigsemstep{e : \tau}{n}
    \monbind
    \lambda e'.\,\sem{e' : \tau}.
  \]
  By the induction hypothesis,
  \[
    \bigsemstep{e : \tau}{n+1}
    =
    \delta_{e : \tau}
    \monbind
    \lambda e'.\,\sem{e' : \tau}.
  \]
  By the left-unit law for Dirac measures,
  \[
    \delta_{e : \tau}
    \monbind
    \lambda e'.\,\sem{e' : \tau}
    =
    \sem{e : \tau}.
  \]
  Since $e \in Val_\tau$, by~\Cref{lem:one-step-value},
  \[
    \sem{e : \tau} = \delta_{e : \tau}.
  \]
  Therefore,
  \[
    \bigsemstep{e : \tau}{n+1}
    =
    \delta_{e : \tau}.
  \]
  Hence the result follows by induction.
\end{proof}

\begin{definition}[Evaluation depth]\label{def:eval-depth}
  The evaluation depth of an expression $\evalDepth \colon Expr_{\tau} \to \mathbb{N}$ is defined as the least number of steps after which the $n$-step distribution is supported only on values, i.e.
  \[
    \evalDepth(e : \tau)
    =
    \min \left\{
      n \in \mathbb{N}
      \mid
      \bigsemstep{e : \tau}{n}(Val_\tau) = 1
    \right\}.
  \]

  Explicitly, this is characterized by~\Cref{fig:eval-depth}, defined recursively on $e : \tau$.
\end{definition}

\begin{figure}
  \small
  \[
  \begin{array}{rcl}
  \evalDepth(v : \tau)
  &=&
  0
  \\[0.75em]

  \evalDepth(\letkw\ x = e_1 : \tau_1\ \inkw\ e_2 : \tau_2)
  &=&
  \begin{cases}
    1 + \evalDepth(e_2[v/x] : \tau_2),
      & \text{if } e_1 = v, \\[0.4em]
    \evalDepth(e_1 : \tau_1)
    + 1
    + \displaystyle\sup_{v \in Val_{\tau_1}}
      \evalDepth(e_2[v/x] : \tau_2),
      & \text{otherwise.}
  \end{cases}
  \\[1.25em]

  \evalDepth(\ifkw\ e_1 : \Bool\ \thenkw\ e_2 : \tau\ \elsekw\ e_3 : \tau)
  &=&
  \begin{cases}
    1 + \evalDepth(e_2 : \tau),
      & \text{if } e_1 = \mathsf{true}, \\[0.4em]
    1 + \evalDepth(e_3 : \tau),
      & \text{if } e_1 = \mathsf{false}, \\[0.4em]
    \evalDepth(e_1 : \Bool)
    + 1
    + \max\{\evalDepth(e_2 : \tau), \evalDepth(e_3 : \tau)\},
      & \text{otherwise.}
  \end{cases}
  \\[1.25em]

  \evalDepth(e_1 + e_2)
  &=&
  \evalDepth(e_1) + \evalDepth(e_2) + 1,
  \\[0.75em]

  \evalDepth(e_1 \times e_2)
  &=&
  \evalDepth(e_1) + \evalDepth(e_2) + 1,
  \\[0.75em]

  \evalDepth(e_1 < e_2)
  &=&
  \evalDepth(e_1) + \evalDepth(e_2) + 1,
  \\[0.75em]

  \evalDepth(\uniform(e_1,e_2))
  &=&
  \evalDepth(e_1) + \evalDepth(e_2) + 1,
  \\[0.75em]

  \evalDepth(\gaussian(e_1,e_2))
  &=&
  \evalDepth(e_1) + \evalDepth(e_2) + 1,
  \\[0.75em]

  \evalDepth(\betafn(e_1,e_2))
  &=&
  \evalDepth(e_1) + \evalDepth(e_2) + 1,
  \\[0.75em]

  \evalDepth(\gammafn(e_1,e_2))
  &=&
  \evalDepth(e_1) + \evalDepth(e_2) + 1,
  \\[0.75em]

  \evalDepth(\exponential(e_1))
  &=&
  \evalDepth(e_1) + 1,
  \\[0.75em]

  \evalDepth(\poisson(e_1))
  &=&
  \evalDepth(e_1) + 1.
  \end{array}
  \]
  \caption{The evaluation depth function $\evalDepth$.}
  \label{fig:eval-depth}
\end{figure}

\begin{example}
  Consider the expression
  \[
    e_1
    =
    (1.0 : \Float{m_1}) + (2.0 : \Float{m_2})
    :
    \Float{m_3}.
  \]
  By the value case of~\Cref{fig:eval-depth},
  \[
    \evalDepth(1.0 : \Float{m_1}) = 0
    \qquad\text{and}\qquad
    \evalDepth(2.0 : \Float{m_2}) = 0.
  \]
  Therefore, by the addition case,
  \[
  \begin{aligned}
    \evalDepth(e_1 : \Float{m_3})
    &=
    \evalDepth(1.0 : \Float{m_1})
    +
    \evalDepth(2.0 : \Float{m_2})
    +
    1 \\
    &= 0 + 0 + 1 \\
    &= 1.
  \end{aligned}
  \]
\end{example}

\begin{example}
  Let
  \[
    e_2 =
    \gaussian(\uniform(0.0,1.0),1.0) : \Float{m}.
  \]
  By the value clause of~\Cref{fig:eval-depth},
  \[
    \evalDepth(0.0 : \Float{m_1}) = 0
    \qquad\text{and}\qquad
    \evalDepth(1.0 : \Float{m_2}) = 0.
  \]
  Therefore, by the uniform clause,
  \[
  \begin{aligned}
    \evalDepth(\uniform(0.0,1.0))
    &=
    \evalDepth(0.0)
    +
    \evalDepth(1.0)
    +
    1 \\
    &= 0 + 0 + 1 \\
    &= 1.
  \end{aligned}
  \]
  Similarly,
  \[
    \evalDepth(1.0 : \Float{m_3}) = 0.
  \]
  Hence, by the Gaussian clause,
  \[
  \begin{aligned}
    \evalDepth(e_2 : \Float{m})
    &=
    \evalDepth(\uniform(0.0,1.0))
    +
    \evalDepth(1.0)
    +
    1 \\
    &= 1 + 0 + 1 \\
    &= 2.
  \end{aligned}
  \]
\end{example}



\begin{lemma}[$n$-step composition]\label{lem:nstep-composition}
  For all $n,k \in \mathbb{N}$,
  \[
    \bigsemstep{e : \tau}{n+k}
    =
    \bigsemstep{e : \tau}{n}
    \monbind
    \lambda e'.\,\bigsemstep{e' : \tau}{k}.
  \]
\end{lemma}

\begin{proof}
  By induction on $k$.

  If $k=0$, then
  \begin{align*}
    \bigsemstep{e : \tau}{n+0} = &\quad \bigsemstep{e : \tau}{n} \\
    = &\quad \bigsemstep{e : \tau}{n}
    \monbind
    \lambda e'.\,\delta_{e' : \tau} \tag{monad right identity}\\
    = &\quad \bigsemstep{e : \tau}{n}
    \monbind
    \lambda e'.\,\bigsemstep{e' : \tau}{0} \tag{definition of $n$-step}.
  \end{align*}

  For the inductive step, assume
  \[
    \bigsemstep{e : \tau}{n+k}
    =
    \bigsemstep{e : \tau}{n}
    \monbind
    \lambda e'.\,\bigsemstep{e' : \tau}{k}.
  \]
  Then, by definition of the $(k+1)$-step semantics and associativity of bind,
  \begin{align*}
    \bigsemstep{e : \tau}{n+(k+1)}
    &=
    \bigsemstep{e : \tau}{(n+k)+1} \\
    &=
    \bigsemstep{e : \tau}{n+k}
    \monbind
    \lambda e'.\,\sem{e' : \tau} \tag{definition of $n$-step}\\
    &=
    \left(
      \bigsemstep{e : \tau}{n}
      \monbind
      \lambda e'.\,\bigsemstep{e' : \tau}{k}
    \right)
    \monbind
    \lambda e''.\,\sem{e'' : \tau} \tag{I.H.}\\
    &=
    \bigsemstep{e : \tau}{n}
    \monbind
    \lambda e'.\,
      \left(
        \bigsemstep{e' : \tau}{k}
        \monbind
        \lambda e''.\,\sem{e'' : \tau}
      \right) \tag{monad assoc.}\\
    &=
    \bigsemstep{e : \tau}{n}
    \monbind
    \lambda e'.\,\bigsemstep{e' : \tau}{k+1} \tag{definition of $n$-step}.
  \end{align*}
  Hence the result follows by induction.
\end{proof}


\begin{lemma}\label{lem:stabilization}
  Let $N = \evalDepth(e : \tau)$. Then for all $n \geq N$,
  \[
    \bigsemstep{e : \tau}{n} = \bigsem{e : \tau}.
  \]
\end{lemma}

\begin{proof}
  Recall by~\Cref{def:eval-depth}, $\bigsemstep{e : \tau}{N}(Val_\tau) = 1$, so $\bigsemstep{e : \tau}{N}$ is supported entirely on values.

  Let $n \geq N$. Then there exists $k \in \mathbb{N}$ such that
  \[
    n = N+k.
  \]

  Then:
  \begin{align*}
    \bigsemstep{e : \tau}{N+k}
    =
    \bigsemstep{e : \tau}{N}
    \monbind
    \lambda e'.\,\bigsemstep{e' : \tau}{k} \tag{by~\Cref{lem:nstep-composition}}\\
    = \bigsemstep{e : \tau}{N}
    \monbind \lambda e'. \delta_{e' : \tau} \tag{by~\Cref{lem:nstep-value}}\\
    = \bigsemstep{e : \tau}{N} \tag{monad right identity}\\
  \end{align*}

  Hence
  \[
    \bigsemstep{e : \tau}{n} = \bigsemstep{e : \tau}{N}
  \]
  for every $n \geq N$.

  Now let $A \subseteq Expr_\tau$ be measurable. Since the sequence
  $\bigsemstep{e : \tau}{m}(A)$ is constant for all $m \geq N$, its limit is the eventual
  constant value:
  \[
    \lim_{m \to \infty} \bigsemstep{e : \tau}{m}(A)
    =
    \bigsemstep{e : \tau}{N}(A).
  \]
  By definition of the big-step semantics,
  \[
    \bigsem{e : \tau}(A)
    =
    \lim_{m \to \infty} \bigsemstep{e : \tau}{m}(A)
    =
    \bigsemstep{e : \tau}{N}(A).
  \]
  Therefore, for every $n \geq N$,
  \[
    \bigsemstep{e : \tau}{n}(A)
    =
    \bigsemstep{e : \tau}{n}(A)
    =
    \bigsemstep{e : \tau}{N}(A)
    =
    \bigsem{e : \tau}(A).
  \]
  Since this holds for every measurable $A$, 
  \[
    \bigsemstep{e : \tau}{n}
    =
    \bigsem{e : \tau}.
  \]
\end{proof}

\begin{lemma}[Expectation for bind]\label{lem:expectation-bind}
  Let $\mu \in \mathcal{D}(X)$ be a measure and let
  $K : X \to \mathcal{D}(\mathbb{R})$ be a Markov kernel. Then
  \[
  \begin{aligned}
    \mathbb{E}[\mu \monbind K]
    &=
    \int y \, d(\mu \monbind K)(y) \\
    &=
    \int
      \left(
        \int y \, dK(a)(y)
      \right)
    \,d\mu(a) \\
    &=
    \int \mathbb{E}[K(a)] \, d\mu(a).
  \end{aligned}
  \]
\end{lemma}

\section{old}
\begin{theorem}[Soundness of determinization]
  Let $e : \tau$ be a well-typed expression. Then:
  \[
    \mathbb{E}[\bigsem{e : \Float{m}}]
    =
    \mathbb{E}[\bigsem{\exptrans{e : \Float{m}} : \Float{m}}]
  \]

  where $m = \G$ or $\E$, and in particular:
  \[
    \bigsem{e : \Float{\G}}
    =
    \bigsem{\exptrans{e : \Float{\G}} : \Float{\G}}.
  \]
\end{theorem}

\begin{proof}
  By induction on $e : \tau$.

  \textbf{Case} $e = \exponential(e_1) : \Float{m}$, where $m$ is either $\E$ or $\G$.

  We want to show:
  \begin{align*}
    \mathbb{E}[\bigsem{\exponential(e_1) : \Float{m}}]
    =
    \mathbb{E}[\bigsem{\exptrans{\exponential(e_1) : \Float{m}} : \Float{m}}] \\
    \iff \\
    \mathbb{E}[\bigsem{\exponential(e_1) : \Float{m}}]
    =
    \mathbb{E}\!\left[
      \bigsem{
        1 : \Float{\E} \;/\; \exptrans{e_1 : \Float{m_1}}
        : \Float{m}
      }
    \right] \tag{by~\Cref{fig:determinization}}.
  \end{align*}

  First, we simplify the LHS expression: 

  Let $d_1 = \evalDepth(e_1 : \Float{m_1})$. By~\Cref{fig:eval-depth},
  \[
    \evalDepth(\exponential(e_1) : \Float{m})
    =
    d_1 + 1.
  \]

  Then:
  \begin{align*}
    &\mathbb{E}[\bigsem{\exponential(e_1) : \Float{m}}]\\
    &=
    \mathbb{E}\!\left[
      \bigsemstep{\exponential(e_1) : \Float{m}}{d_1+1}
    \right] \tag{by~\Cref{lem:stabilization}}\\
    &=
    \mathbb{E}\!\left[
      \bigsemstep{\exponential(e_1) : \Float{m}}{d_1}
      \monbind
      \lambda e'.\,\sem{e' : \Float{m}}
    \right]\\
    &=
    \mathbb{E}\!\left[
      \bigsemstep{e_1 : \Float{m_1}}{d_1}
      \monbind
      \lambda a.\,
        \delta_{\exponential(a) : \Float{m}}
      \monbind
      \lambda e'.\,\sem{e' : \Float{m}}
    \right] \tag{???idk}\\
    &=
    \mathbb{E}\!\left[\bigsemstep{e_1 : \Float{m_1}}{d_1}
    \monbind
    \lambda a.\,
      \sem{\exponential(a) : \Float{m}} \right] \\
    &=
    \mathbb{E}\!\left[\bigsem{e_1 : \Float{m_1}}
    \monbind
    \lambda a.\,
      \sem{\exponential(a) : \Float{m}}\right] \tag{by~\Cref{lem:stabilization}}\\
  \intertext{Since $a$ is a value, the small-step semantics for $\exponential$ gives
  $\sem{\exponential(a) : \Float{m}}
  =
  \mathsf{Exponential}(a)
  \monbind
  \lambda x.\,\delta_{x : \Float{m}}$. Hence,}
    &=
    \mathbb{E}\!\left[\bigsem{e_1 : \Float{m_1}}
    \monbind
    \lambda a.\,
      \left(
        \mathsf{Exponential}(a)
        \monbind
        \lambda x.\,\delta_{x : \Float{m}}
      \right)\right] \tag{by the small-step semantics for $\exponential$}\\
    &=
    \int
      \mathbb{E}\!\left[
        \mathsf{Exponential}(a)
        \monbind
        \lambda x.\,\delta_{x : \Float{m}}
      \right]
    \,d\bigsem{e_1 : \Float{m_1}}(a) \tag{by~\Cref{lem:expectation-bind}}\\
    &=
    \int
      \left(
      \int
        \mathbb{E}\!\left[
          \delta_{x : \Float{m}}
        \right]
      \,d\mathsf{Exponential}(a)(x)
      \right)
    \,d\bigsem{e_1 : \Float{m_1}}(a) \\
    &=
    \int
      \left(
      \int
        x
      \,d\mathsf{Exponential}(a)(x)
      \right)
    \,d\bigsem{e_1 : \Float{m_1}}(a) \\
    &=
    \int
      \mathbb{E}\!\left[
        \mathsf{Exponential}(a)
      \right]
    \,d\bigsem{e_1 : \Float{m_1}}(a) \\
    &=
    \int
      \frac{1}{a}
    \,d\bigsem{e_1 : \Float{m_1}}(a).
  \end{align*}

  Now simplify the LHS expression:

  By~\Cref{fig:eval-depth},
  \[
    \evalDepth(1 : \Float{\E} \;/\; \exptrans{e_1 : \Float{m_1}} : \Float{m})
    =
    \evalDepth(1 : \Float{\E}) + \evalDepth(\exptrans{e_1 : \Float{m_1}} : \Float{m_1}) + 1.
  \]

  Since $1 : \Float{\E}$ is a value,
  \[
    \evalDepth(1 : \Float{\E}) = 0.
  \]

  Let
  \[
    d_1' = \evalDepth(\exptrans{e_1 : \Float{m_1}} : \Float{m_1}).
  \]

  Then:
  \begin{align*}
    &\mathbb{E}\!\left[
      \bigsem{
        1 : \Float{\E} \;/\; \exptrans{e_1 : \Float{m_1}}
        : \Float{m}
      }
    \right]\\
    &=
    \mathbb{E}\!\left[
      \bigsemstep{
        1 : \Float{\E} \;/\; \exptrans{e_1 : \Float{m_1}}
        : \Float{m}
      }{d_1'+1}
    \right] \tag{by~\Cref{lem:stabilization}}\\
    &=
    \mathbb{E}\!\left[
    \bigsemstep{\exptrans{e_1 : \Float{m_1}} : \Float{m_1}}{d_1'}
    \monbind
    \lambda a.\,
      \sem{1 : \Float{\E} \;/\; a : \Float{m}} \right] \tag{???}\\
    &=
    \mathbb{E}\!\left[
    \bigsem{\exptrans{e_1 : \Float{m_1}} : \Float{m_1}}
    \monbind
    \lambda a.\,
      \sem{1 : \Float{\E} \;/\; a : \Float{m}}\right] \tag{by~\Cref{lem:stabilization}}\\
  \intertext{Since $a$ is a value, the small-step semantics for division gives
  $\sem{1 : \Float{\E} \;/\; a : \Float{m}}
  =
  \delta_{(1/a) : \Float{m}}$. Hence,}
    &=
    \mathbb{E}\!\left[
    \bigsem{\exptrans{e_1 : \Float{m_1}} : \Float{m_1}}
    \monbind
    \lambda a.\,
      \delta_{(1/a) : \Float{m}}\right] \tag{by the small-step semantics for division}\\
    &=
    \int
      \mathbb{E}\!\left[
        \delta_{(1/a) : \Float{m}}
      \right]
    \,d\bigsem{\exptrans{e_1 : \Float{m_1}} : \Float{m_1}}(a) \tag{by~\Cref{lem:expectation-bind}}\\
    &=
    \int
      \left(
      \int
        y
      \,d\delta_{(1/a) : \Float{m}}(y)
      \right)
    \,d\bigsem{\exptrans{e_1 : \Float{m_1}} : \Float{m_1}}(a) \\
    &=
    \int
      \frac{1}{a}
    \,d\bigsem{\exptrans{e_1 : \Float{m_1}} : \Float{m_1}}(a) \\
    &=
    \int
      \frac{1}{a}
    \,d\bigsem{\exptrans{e_1 : \Float{m_1}}}(a).
  \end{align*}

  Thus, this reduces to showing:
  \[
    \int
      \frac{1}{a}
    \,d\bigsem{e_1 : \Float{m_1}}(a)
    =
    \int
      \frac{1}{a}
    \,d\bigsem{\exptrans{e_1 : \Float{m_1}}}(a).
  \]
  Because $m_1 = \G$, $\bigsem{e_1 : \Float{\G}} = \bigsem{\exptrans{e_1 : \Float{\G}} : \Float{\G}}$ by the I.H.

  Therefore,
  \[
    \mathbb{E}[\bigsem{\exponential(e_1) : \Float{m}}]
    =
    \mathbb{E}\!\left[
      \bigsem{
        1 : \Float{\E} \;/\; \exptrans{e_1 : \Float{m_1}}
        : \Float{m}
      }
    \right].
  \]
\end{proof}


\newpage
\section{old}
\subsection{Transformation}

\begin{definition}[Expectation transformation]
  Define $\exptrans{e : \mathsf{float}}$ inductively on the structure of $e$.

  *** This will be dependent on the types.
\end{definition}

\begin{lemma}
  Type preservation of the expectation transformer.
\end{lemma}


\subsection{Expectations}

\begin{lemma}[Expectation is well-defined]
  For any sub-probability measure $\mu \in \mathcal{D}(Vals_{\tau})$, the expectation
  \[
    \mathbb{E}[\mu] := \int x \, d\mu(x)
  \]
  is a well-defined real number.

  **Specifically, let us take $\tau = \mathbf{float}$.
\end{lemma}

\begin{lemma}[Expectation of a Dirac]
  For any $c \in Vals_{\tau}$:
  \[
    \mathbb{E}[\delta_c] = c
  \]
\end{lemma}

\begin{lemma}[Linearity of expectation]
  For sub-probability measures $\mu_1, \mu_2 \in \mathcal{D}(Vals_{\tau})$ and $a, b \in \mathbb{R}$:
  \[
    \mathbb{E}[a \cdot \mu_1 + b \cdot \mu_2] = a \cdot \mathbb{E}[\mu_1] + b \cdot \mathbb{E}[\mu_2]
  \]
\end{lemma}

\begin{lemma}
  For $\mu \in \mathcal{D}(Expr_\tau)$ and measurable $K : Expr_\tau \to \mathcal{D}(Vals_{\tau})$:
  \[
    \mathbb{E}[\mu \monbind K] = \int \mathbb{E}[K(x)] \, d\mu(x)
  \]
\end{lemma}

\begin{lemma}[Expectation of a program is well-defined]
  For any well-typed expression $e : \mathsf{float}$, the expectation
  \[
    \mathbb{E}[\bigsem{e : \mathsf{float}}] 
  \]
  is a well-defined real number.
\end{lemma}

\subsection{Affine condition}

\begin{definition}[Affine continuation]
  A measurable function $K : \mathbb{R} \to \mathcal{D}(\mathbb{R})$ is called \emph{expectation-affine} if the map
  \[
    x \mapsto \mathbb{E}[K(x)]
  \]
  is an affine function of $x$, i.e.\ there exist $a, b \in \mathbb{R}$ such that $\mathbb{E}[K(x)] = ax + b$ for all $x$.
\end{definition}

\begin{itemize}
  \item The continuation $K(x) = \delta_{ax+b}$ is expectation-affine. 
  \item If $K_1 : \mathbb{R} \to \mathcal{D}(\mathbb{R})$ is expectation-affine and $K_2 : \mathbb{R} \to \mathcal{D}(\mathbb{R})$ is expectation-affine, then their monadic composition
  \[
    x \mapsto K_1(x) \monbind K_2
  \]
  is expectation-affine.
  \item The continuation $K(x) = \mathcal{N}(ax + b, \sigma^2)$ is expectation-affine, with $\mathbb{E}[K(x)] = ax + b$.
  \item The continuation $K(x) = \mathrm{Uniform}(ax + b,\, cx + d)$ is expectation-affine, with
  \[
    \mathbb{E}[K(x)] = \frac{(a+c)}{2} x + \frac{(b+d)}{2}.
  \]
\end{itemize}

\begin{definition}[Exepectation equivalence]
  Define $K \colon Val_\tau \to \mathcal{D}(\mathbb{R})$ respects expectation equivalence if for all $\mu_1, \mu_2$ such that $\mu_1 \equiv_\tau \mu_2$, then $\mathbb{E}[\mu_1 \monbind K] = \mathbb{E}[\mu_2 \monbind K]$.
\end{definition}

\begin{definition}[Admissible continuations]
  For each type $\tau$, an observation continuation is a Markov kernel
  \[
    K : Val_\tau \to \mathcal{D}(\mathbb{R}).
  \]
  Write
  \[
    \overline{K}(v) = \mathbb{E}[K(v)].
  \]
  We define a type-indexed predicate $Adm_\tau(K)$ by cases on $\tau$.

  \begin{itemize}
    \item $\tau = \Unit$. Every continuation is admissible.

    \item $\tau = \Bool$. Every continuation is admissible.

    \item $\tau = \Float{\G}$. Every continuation is admissible.

    \item $\tau = \Float{\E}$. The continuation is admissible iff
    $\overline{K}$ is an affine continuation:
    \[
      \exists a,b \in \mathbb{R}.\ \forall x \in \mathbb{R},\quad
      \overline{K}(x) = ax + b.
    \]

    \item $\tau = \tau_1 \times \tau_2$. The continuation $K$ is
    admissible iff it is admissible in each component separately:
    \[
      \forall v_2 \in Val_{\tau_2},\quad
      Adm_{\tau_1}(\lambda v_1.\,K(v_1,v_2))
    \]
    and
    \[
      \forall v_1 \in Val_{\tau_1},\quad
      Adm_{\tau_2}(\lambda v_2.\,K(v_1,v_2)).
    \]

    \item $\tau = \tau_1 + \tau_2$. The continuation $K$ is admissible iff
    both branch continuations are admissible:
    \[
      Adm_{\tau_1}(\lambda v.\,K(\mathsf{inl}\,v))
    \]
    and
    \[
      Adm_{\tau_2}(\lambda v.\,K(\mathsf{inr}\,v)).
    \]

    \item $\tau = Option{\tau'}$. The continuation $K$ is admissible iff
    the $\mathsf{some}$-branch continuation is admissible:
    \[
      Adm_{\tau'}(\lambda v.\,K(\mathsf{some}\,v)).
    \]
    The $\mathsf{none}$ branch is a single point and imposes no condition.

    \item $\tau = \List{\tau'}$. The continuation $K$ is admissible iff,
    for every list context $C[-]$ with one hole of type $\tau'$, the induced
    element continuation is admissible:
    \[
      Adm_{\tau'}(\lambda v.\,K(C[v])).
    \]
    Equivalently, for every list length $n$, every position $i$, and every
    fixed choice of all other elements, the function obtained by varying only
    the $i$-th element is admissible at type $\tau'$.

    \item $\tau = \tau_1 \to \tau_2$. The continuation $K$ is admissible
    extensionally: it may observe a function only by applying it to admissible
    arguments and then observing the result admissibly. Concretely, $K$ is
    admissible if it belongs to the least class of continuations closed under
    constants, affine combinations, and the rule:
    if $A$ is an admissible argument-producing computation of type
    $\tau_1$, and $L$ is admissible at type $\tau_2$, then
    \[
      f \mapsto \bigsem{f\,A : \tau_2} \monbind L
    \]
    is admissible at type $\tau_1 \to \tau_2$.
  \end{itemize}
\end{definition}


\begin{lemma}
  For any $\mu \in \mathcal{D}(Expr_\mathsf{float})$ and any expectation-affine continuation $K$:
  \[
    \mathbb{E}[\mu \monbind K] = \mathbb{E}[\delta_{\mathbb{E}[\mu]} \monbind K]
  \]
\end{lemma}

\begin{proof}
  \begin{align*}
    &\mathbb{E}[\mu \monbind K] \\
    = \quad&\int \mathbb{E}[K(x)] \, d\mu(x) \tag{Lemma 5.12} \\
    = \quad&\int (ax + b) \, d\mu(x) \tag{Definition 5.14} \\
    = \quad&a \cdot \mathbb{E}[\mu] + b \tag{Linearity of expectation} \\
    = \quad&\mathbb{E}[K(\mathbb{E}[\mu])] \tag{Definition 5.14} \\
    = \quad&\mathbb{E}[\delta_{\mathbb{E}[\mu]} \monbind K] 
  \end{align*}
\end{proof}

\begin{theorem}[Soundness under affine continuations]
  Let $e : \tau$ be well typed. Let
  $K : Val_\tau \to \mathcal{D}(\mathbb{R})$
  be an admissible continuation. Then
  \[
    \mathbb{E}[\bigsem{e : \tau} \monbind K]
    =
    \mathbb{E}[\bigsem{\exptrans{e : \tau} : \tau} \monbind K].
  \]
\end{theorem}

\begin{theorem}[Soundness of determinization]
  Let $e : \tau$ be a well-typed expression. Then:
  \[
    \mathbb{E}[\bigsem{e : \Float{m}}]
    =
    \mathbb{E}[\bigsem{\exptrans{e : \Float{m}} : \Float{m}}]
  \]

  where $m = \G$ or $\E$, and in particular:
  \[
    \bigsem{e : \Float{\G}}
    =
    \bigsem{\exptrans{e : \Float{\G}} : \Float{\G}}.
  \]
\end{theorem}

\begin{proof}
  Let $\tau = \Float{m}$ and $K(x) = \delta_x$. 
\end{proof}
